\documentclass[../../main2.tex]{subfiles}

\begin{document}

	% ======================================================================
	% 骨架文件：小节标题与追问链为终版；正文由后续会话填充。
	% 扩增模式标注：本章 = 无压扩增（传播不消耗原体）。
	% 原书材料接入点：
	%   - 原书 part3 ch2（文化与民间制度）→ 本章主体
	%   - 原书 part1 ch3（语言与信息）→ §9.3：语言 = 模因的最高频压缩层
	%     （从「核心逻辑」降级至此，是有意改动）
	% ======================================================================

	\chapter{通道一：模因（无压扩增）}
	\label{ch:meme}

	\begin{motivation}
		上一章遗留的问题：生成树的第一个节点——不加任何约束的传播长什么样？什么信息会自己传播？\\
		\textbf{核心动机}：定义模因（meme）并论证其无压属性；确立其关键指标是传播强度（intensity）而非方向（direction）；给出模因通道的贡献分解。\\
		\textbf{扩增模式}：无压——复制不消耗原体，上限在接收方。
	\end{motivation}

	\begin{logicmap}
	\begin{tikzpicture}[node distance=0.9cm, every node/.style={draw, rectangle, rounded corners, align=center}]
	\node (q) {什么信息会自己传播？};
	\node (d) [below=of q] {模因定义与判据（§9.1）};
	\node (s) [below=of d] {传播强度 > 方向正确（§9.2）};
	\node (l) [below=of s] {语言 = 压缩层（§9.3）};
	\node (c) [below=of l] {模因的贡献分解（§9.4）};
	\node (n) [below=of c] {传播要花钱呢？（抛给第10章）};
	\draw[-{Stealth}] (q) -- (d);
	\draw[-{Stealth}] (d) -- (s);
	\draw[-{Stealth}] (s) -- (l);
	\draw[-{Stealth}] (l) -- (c);
	\draw[-{Stealth}] (c) -- (n);
	\end{tikzpicture}
	\end{logicmap}

	% ==================================================================
	\section{模因的定义与判据}
	模因（meme）是文化通道的基本单元：任何可以被复制、传播、且复制后原体保留的信息模式。一个笑话、一段旋律、一种烹饪方法、一套道德规范——都是模因。

	判据有三条，缺一不可：(1) 可复制性——存在从载体 A 到载体 B 的传递机制；(2) 原体保留——传递后源载体不失去该信息（$A \to A+B$，不是 $A \to B$）；(3) 选择压力——不同变体在传播成功率上有差异（有的笑话被转发一万次，有的只被转了一次）。

	第三条容易忽略但至关重要：没有选择压力的复制只是噪声。模因和随机噪声的区别在于，模因的种群分布是非均匀的——某些变体系统性地产出更多副本。这个「系统性偏差」就是文化演化的驱动力，和生物进化里的自然选择同构，只是作用对象从基因换成了信息模式。
	% ==================================================================
	\section{传播强度决定一切}
	模因通道的核心变量是传播强度：一个模因在单位时间内能产生多少副本。

	传播强度取决于三个因素：(1) 保真度——复制时丢失多少信息（口耳相传 vs 文字记录 vs 数字存储）；(2) 载体密度——潜在接收者的数量和接触频率（朋友圈 vs 村口大树下）；(3) 选择偏好——接收者主动筛选的倾向（算法推荐 vs 随机广播）。

	无压扩增意味着没有总量上限：一个模因可以理论上复制到无穷多个载体。但实际中传播强度有衰减——每一代复制都引入噪声，载体注意力有限，竞争模因争夺同一批接收者。所以真实的文化演化不是指数爆炸，而是带饱和的 logistic 增长：早期快速扩散，后期趋于平台。

	\begin{readingnote}
		为什么有些歌曲火了一辈子，有些只火了一个月？保真度相同（都是数字音频），载体密度相同（都在同一个平台上），区别在选择偏好：前者击中了跨代际的共同情感结构，后者只击中了某个亚文化的即时情绪。传播强度的差异不在技术层，在内容层。
	\end{readingnote}
	% ==================================================================
	\section{语言作为压缩层}
	语言是模因传播的压缩层。

	没有语言之前，文化传递靠示范和模仿：你看到别人做，你学着做。保真度极低——动作的微妙差异在观察-复制过程中大量丢失。语言把连续的动作空间离散成有限的符号集合（词汇），把高维的行为模式压缩成低维的序列（句子）。压缩带来了两个好处：保真度提升（「把盐放在左边第三个抽屉」比示范精确得多），传播速度提升（一句话一秒钟，示范要五分钟）。

	代价是信息损失：压缩必然有损。语言无法完整编码所有感官细节、情感色调、情境上下文。这就是为什么「如临其境」永远是修辞而非事实——语言是压缩编码，解压后的重建不等于原始信号。第12章的法律变长编码和这里的语言压缩是同构的：都是高频模式用短码字、低频模式用长码字。
	% ==================================================================
	\section{模因的贡献分解}
	把模因通道的贡献写进分解公式：$C_{\text{模因}}$ 衡量的是「纯粹由信息传播驱动的行为变化」在总行为中的份额。

	测量方法：对比「知道某信息」和「不知道某信息」两种状态下行为的差异，控制资源约束和身份关系不变。例如：一条健康科普视频发布前后，某社区的运动率变化——如果资源（健身房可及性）和身份（社会规范）没变，差值就是 $C_{\text{模因}}$。

	典型量级：在信息密集的社会变迁中（技术扩散、观念转变），$C_{\text{模因}}$ 往往占主导（50-80%）；在资源瓶颈型变迁中（饥荒、战争），$C_{\text{经济}}$ 占主导。两者不是此消彼长——可以同时显著，但归因时必须分开算。
	% ==================================================================
	\section{本章小结与衔接}
	模因 = 无压复制 + 选择压力。它是四条通道的根节点：一切社会行为都涉及信息传播，经济、政治、法律都是在这个基底层上加了额外约束的子集。传播强度决定扩散速度，语言是压缩层（提高保真度和速度，代价是信息损失），贡献分解通过控制变量法隔离纯信息效应。

	\paragraph*{逻辑衔接}
	第10章「通道二：经济（模因 + 稀缺约束）」：当传播必须付费时，守恒律登场。
\end{document}
